\section{Related Works}

\paragraph{Multi-agent Debate}

Proposed by MAD \cite{liang2023encouraging} and SoM \cite{du2023improving}, Multi-Agent Debate has demonstrated significant performance improvements across diverse reasoning benchmarks. ChatEval \cite{chan2023chateval} utilizes this interaction mechanism for autonomous evaluation. Some research focuses on optimizing communication topologies \cite{li2024improving} to reduce token consumption, while other work integrates confidence scores \cite{lin2025enhancing} to weigh agent contributions. Regarding heterogeneity, Heter-MAD \cite{zhang2025if} investigates the impact of diversity, and DMAD  \cite{liu2025breaking} employs diverse prompting strategies to foster divergent reasoning.

\paragraph{Tool-Augmented Reasoning and Verification}

Recent works augment LLMs with external utilities to overcome closed-loop limitations. Toolformer \cite{schick2023toolformer} and ART \cite{paranjape2023art} enable models to autonomously invoke tools for multi-step reasoning, while ToolLLM \cite{qin2023toolllm} scales this capability to master real-world APIs. In mathematical domains, Program of Thoughts (PoT) \cite{chen2022program} disentangles computation from reasoning, and Code-based Self-Verification \cite{zhou2023solving} leverages code execution to verify solutions. Focusing on robustness, CRITIC \cite{gou2023critic} interacts with tools to validate outputs, and VerifiAgent \cite{han2025verifiagent} proposes a unified verification framework for reasoning tasks.

\paragraph{Multi-path reasoning}

To mitigate the instability of single-chain reasoning, researchers have developed multi-path strategies that explore diverse solution spaces. Self-Consistency (CoT-SC) \cite{wang2022self} pioneered this direction by sampling multiple reasoning paths and applying majority voting to filter out random errors. Moving beyond independent sampling, Tree of Thoughts (ToT) \cite{yao2023tree} structures reasoning as a tree search, enabling lookahead and backtracking via algorithms like BFS or DFS. Graph of Thoughts (GoT) \cite{besta2024graph} further generalizes this into a graph topology, allowing for the aggregation and recurrence of thoughts.
